
\section{\BookRecordIndex{T0330}椭圆自举}\label{sec:B-4}

本节将证明以下两个定理：其一关于具有全实边界条件的 \BookRecordIndex{T0003}$J$-全纯曲线的光滑性，
其二关于导数的 $L^p$ 范数（$p>2$）具有一致上界的序列的\BookRecordIndex{T0192}紧性。
设 $M$ 是光滑的 $2n$ 维流形，$L\subset M$ 是闭的 $n$ 维子流形。
以 $\mathcal J^\ell(M,L)$ 表示 $M$ 上使 $L$ 为\BookRecordIndex{T0290}全实子流形的 $C^\ell$ 类\BookRecordIndex{T0088}殆复结构 $J$ 的集合
（即 $TL\cap J(TL)=\{0\}$）。当 $\ell=\infty$ 时，记
$\mathcal J(M,L):=\mathcal J^\infty(M,L)$。
设 $\Sigma$ 是带边界的定向二维流形，以 $\mathcal J(\Sigma)$ 表示 $\Sigma$ 上复结构的集合。

\begin{theorem}[\BookRecordIndex{T0408}正则性]\label{thm:B-4-1}
固定整数 $\ell\ge2$ 和数 $p>2$。设 $j\in\mathcal J(\Sigma)$、
$J\in\mathcal J^\ell(M,L)$，并设 $u:\Sigma\to M$ 是 $W^{1,p}$ 类映射，满足
$du\circ j=J\circ du$ 及 $u(\partial\Sigma)\subset L$。
则 $u$ 为 $W^{\ell,p}$ 类。若 $\ell=\infty$，则 $u$ 光滑。
\sourcefootnote{原书未要求 $\Sigma$ 紧致，却给出全局 $W^{\ell,p}$ 结论；
下文证明实际得到的是 $W^{\ell,p}_{\mathrm{loc}}$。在非紧定义域上，
局部\BookRecordIndex{T0408}正则性不自动给出全局 索伯列夫范数有限；此处保留原断言。}
\end{theorem}

\begin{theorem}[\BookRecordIndex{T0192}紧性]\label{thm:B-4-2}
固定整数 $\ell\ge2$ 和数 $p>2$。设 $J_\nu\in\mathcal J^\ell(M,L)$
是一列\BookRecordIndex{T0088}殆复结构，在 $C^\ell$ 拓扑下收敛到 $J\in\mathcal J^\ell(M,L)$；
设 $j_\nu$ 是 $\Sigma$ 上一列复结构，在 $C^\infty$ 拓扑下收敛到 $j$。
设 $U_\nu\subset\Sigma$ 是并集为 $\Sigma$ 的递增开集列，
$u_\nu:U_\nu\to M$ 是一列 $W^{1,p}$ 类的 $(j_\nu,J_\nu)$-全纯曲线，且
\[
  u_\nu(U_\nu\cap\partial\Sigma)\subset L.
\]
假设对每个具有光滑边界的紧集 $Q\subset\Sigma$，存在紧集 $K\subset M$
及常数 $c>0$，使得当 $\nu$ 充分大时
\[
  \|du_\nu\|_{L^p(Q)}\le c,\qquad u_\nu(Q)\subset K.
\]
则 $u_\nu$ 有一子序列，在 $\Sigma$ 的每个紧子集上以 $C^{\ell-1}$ 拓扑收敛。
\end{theorem}

\begin{remark}\label{rem:B-4-3}
若 $\partial\Sigma=\varnothing$，则只需 $\ell\ge1$；
在定理~\ref{thm:B-4-1} 中得到 $u\in W^{\ell+1,p}_{\mathrm{loc}}$，
在定理~\ref{thm:B-4-2} 中得到一个 $C^\ell$ 收敛的子序列。
此外，略为细致的讨论表明，采用 $C^{\ell-1,1}$ 类\BookRecordIndex{T0088}殆复结构已足够。
特别地，对内部\BookRecordIndex{T0408}正则性而言，可以采用 Lipschitz 连续的\BookRecordIndex{T0088}殆复结构。
\end{remark}

当 $J$ 可积且 $\Sigma$ 为闭曲面时，这两个定理显然成立，因为在全纯坐标中，
\BookRecordIndex{T0003}$J$-全纯曲线的每个分量都是\BookRecordIndex{T0104}调和函数。特别地，\BookRecordIndex{T0192}紧性定理由\BookRecordIndex{T0104}调和函数的\BookRecordIndex{T0267}平均值性质\BookRecordIndex{T0326}推出。
一般情形则困难得多。

两个定理的证明的主要思路，是以线性方程 $\partial_su+J'\partial_tu=0$
代替非线性方程 $\partial_su+J(u)\partial_tu=0$；这里
$J'(z)=J(u(z))$ 是定义在上半平面某个开子集 $\Omega\subset\mathbb H$
上的矩阵值函数。这需要在 $M$ 上选取适当坐标，使边界条件也成为线性的。
注意 $J'$ 与 $u$ 具有相同的光滑性，这便导向下文的\BookRecordIndex{T0330}椭圆自举论证。
更具体地，设 $J\in W^{1,p}(\Omega,\mathbb R^{2n\times2n})$ 满足
$J^2=-\mathbf 1$ 以及
\begin{equation}\label{eq:B-4-1}
  J(s,0)=J_0:=\begin{pmatrix}0&-\mathbf 1\\ \mathbf 1&0\end{pmatrix},
\end{equation}
其中 $J_0$ 是 $\mathbb R^{2n}$ 上的标准复结构。
设 $u\in W^{1,p}(\Omega,\mathbb R^{2n})$ 是边值问题
\[
  \partial_su+J\partial_tu=0,\qquad u(s,0)\in\mathbb R^n\times\{0\}
\]
的解。则导数 $v:=\partial_su$ 满足非齐次方程
\[
  \partial_sv+J\partial_tv=\eta,\qquad
  v(s,0)\in\mathbb R^n\times\{0\},\qquad
  \eta:=-(\partial_sJ)\partial_tu.
\]
从\BookRecordIndex{T0408}正则性定理中我们至多能期望：若 $\eta\in L^r(\Omega)$，其中 $r>1$，
则 $v\in W^{1,r}(\Omega)$。我们希望取 $r=p$，但假设只说明
$\partial_sJ,\partial_tu\in L^p$，故它们的乘积一般仅属于 $L^{p/2}$。
为解决这一问题，我们在保持 $\partial_sJ$ 的\BookRecordIndex{T0408}正则性不变的同时提高 $u$ 的\BookRecordIndex{T0408}正则性。
假设已证明 $\partial_tu\in L^q$，其中 $q>p/(p-1)$。
于是 $\eta\in L^r$，其中 $1/p+1/q=1/r$。由于 $r>1$，便有
$v\in W^{1,r}$，从而 $u\in W^{2,r}$。若 $1<r<2$，
则由定理~\ref{thm:B-1-12}，$u\in W^{1,q'}$，其中
$q':=2r/(2-r)>q$。因此\BookRecordIndex{T0408}正则性有所提高。现在重复此过程。
一旦得到大于 $2$ 的 $r$，定理~\ref{thm:B-1-11} 就表明
$\partial_tu$ 连续，因而 $\eta\in L^p$。
这一迭代是\BookRecordIndex{T0330}椭圆自举论证的一部分；它需要下面两个引理。

\begin{lemma}\label{lem:B-4-4}
给定 $p>2$，存在有限序列 $q_0<q_1<\cdots<q_m$，使得对
$j=0,\ldots,m-1$，有
\[
  \frac p{p-1}<q_0\le p,\qquad
  q_{m-1}<\frac{2p}{p-2}<q_m,
\]
以及
\[
  q_{j+1}=\frac{2r_j}{2-r_j},\qquad r_j:=\frac{pq_j}{p+q_j}.
\]
下文出现的 $r_m$ 也按同一公式定义。
\end{lemma}
\begin{proof}
考虑映射 $h:(p/(p-1),2p/(p-2))\to(2,\infty)$，其定义为
\[
  h(q):=\frac{2pq}{2p+2q-pq}=\frac{2r}{2-r},\qquad
  r:=\frac{pq}{p+q}<2.
\]
条件 $r<2$ 等价于 $q<2p/(p-2)$。映射 $h$ 是单调递增的微分同胚，
且 $h(q)>q$。取序列 $q_j$ 满足 $q_{j+1}=h(q_j)$ 即可。
\end{proof}

\begin{lemma}\label{lem:B-4-5}
设 $p>2$ 且 $1<r\le p$。则存在常数 $c=c(p,r)>0$，使得：
若 $f\in W^{1,p}(\mathbb R^2)$、$g\in W^{1,r}(\mathbb R^2)$，
则 $fg\in W^{1,r}(\mathbb R^2)$，且
\[
  \|fg\|_{W^{1,r}}\le c\|f\|_{W^{1,p}}\|g\|_{W^{1,r}}.
\]
\end{lemma}
\begin{proof}
考察 $d(fg)=(df)g+f(dg)$ 的 $L^r$ 范数。
项 $f(dg)$ 可由 $f$ 的上确界范数与 $g$ 的 $W^{1,r}$ 范数估计。
对项 $(df)g$，有
\[
  \|(df)g\|_{L^r}\le\|df\|_{L^p}\|g\|_{L^q},\qquad
  1/p+1/q=1/r.
\]
若 $r<2$，由于 $p>2$，有 $q=pr/(p-r)<2r/(2-r)$，
故 $g$ 的 $L^q$ 范数可由其 $W^{1,r}$ 范数估计。
当 $r\ge2$ 时，后一估计显然成立。
\end{proof}

接下来证明线性\BookRecordIndex{T0209}柯西–黎曼方程的解的局部估计。
设 $\mathbb H:=\{s+it\in\mathbb C\mid t\ge0\}$ 为闭上半平面，
$\Omega\subset\mathbb H$ 为开集。注意 $\Omega$ 未必是 $\mathbb C$ 中的开集。
以 $C_0^\infty(\Omega,\mathbb R^{2n})$ 表示支集为 $\Omega$ 中紧集的
光滑函数 $\phi:\Omega\to\mathbb R^{2n}$ 的空间。
这些函数在边界部分 $\Omega\cap\mathbb R$ 上不必为零。
同前，以 $W^{k,p}_{\mathrm{loc}}(\Omega)$ 表示在 $\Omega$ 的每个紧子集上
限制为 $W^{k,p}$ 类的函数空间；但在这里，此类函数也须在
$\Omega\cap\mathbb R$ 附近具有 $W^{k,p}$ \BookRecordIndex{T0408}正则性。
设 $u:\Omega\to\mathbb R^{2n}$ 是线性边值问题
\begin{equation}\label{eq:B-4-2}
  \partial_su+J\partial_tu=\eta,\qquad
  u(\Omega\cap\mathbb R)\subset\mathbb R^n\times\{0\}
\end{equation}
的解，其中 $\eta:\Omega\to\mathbb R^{2n}$、
$J:\Omega\to\mathbb R^{2n\times2n}$ 满足 $J^2=-\mathbf 1$ 及~\eqref{eq:B-4-1}。
先在对\BookRecordIndex{T0088}殆复结构 $J$ 和函数 $\eta$ 尽可能弱的\BookRecordIndex{T0408}正则性假设下证明 $u$ 的\BookRecordIndex{T0408}正则性。
首先假设 $J$ 为 $W^{1,p}$ 类。若 $u$ 是~\eqref{eq:B-4-2} 的解，
则分部积分给出
\begin{equation}\label{eq:B-4-3}
  \phi(\Omega\cap\mathbb R)\subset\mathbb R^n\times\{0\}
  \quad\Longrightarrow\quad
  \int_\Omega\langle\partial_s\phi+J^T\partial_t\phi,u\rangle
  =-\int_\Omega\langle\phi,\eta+(\partial_tJ)u\rangle
\end{equation}
对每个检验函数 $\phi\in C_0^\infty(\Omega,\mathbb R^{2n})$ 成立，
其中 $J^T$ 表示 $J$ 的转置。下一个引理断言其逆命题。

\begin{lemma}\label{lem:B-4-6}
设 $\Omega'\subset\Omega\subset\mathbb H$ 为开集，满足
$\overline{\Omega'}\subset\Omega$\sourcefootnote{原书仅要求
$\overline{\Omega'}\subset\Omega$，没有要求 $\Omega'$ 有界；若 $\Omega'$ 非有界，
这一包含关系并不保证统一的内缩距离，不能仅凭逐紧局部估计得到下文
对整个 $\Omega'$ 使用同一常数的估计。证明实际先在有界 Lipschitz 子域上估计；
命题~\ref{prop:B-4-9} 沿用相同开集条件，也有此疑点。此处保留原假设。}；
设 $p,q,r$ 为正实数（亦允许正无穷），满足
\[
  2<p,\qquad 1<r<\infty,\qquad 1/p+1/q=1/r.
\]
则对每个常数 $c_0>0$，存在常数 $c>0$，使得如下结论成立。
假设 $J\in W^{1,p}(\Omega,\mathbb R^{2n\times2n})$ 满足
$J^2=-\mathbf 1$、\eqref{eq:B-4-1}，且
\[
  \|J\|_{W^{1,p}(\Omega)}\le c_0.
\]
则有以下结论。
\begin{enumerate}[label=\textup{(\roman*)}]
\item 若 $u\in L^q_{\mathrm{loc}}(\Omega,\mathbb R^{2n})$ 和
$\eta\in L^r_{\mathrm{loc}}(\Omega,\mathbb R^{2n})$ 使~\eqref{eq:B-4-3} 成立，
则 $u\in W^{1,r}_{\mathrm{loc}}(\Omega,\mathbb R^{2n})$，
且 $u$ 几乎处处满足~\eqref{eq:B-4-2}。
\item 若 $u\in W^{1,r}_{\mathrm{loc}}(\Omega)$ 满足边界条件
$u(\Omega\cap\mathbb R)\subset\mathbb R^n\times\{0\}$，则
\[
  \|u\|_{W^{1,r}(\Omega')}\le
  c\bigl(\|\partial_su+J\partial_tu\|_{L^r(\Omega)}
  +\|u\|_{L^q(\Omega)}\bigr).
\]
\end{enumerate}
\end{lemma}

\begin{remark}\label{rem:B-4-7}
注意引理~\ref{lem:B-4-6} 的假设包括 $q=\infty$ 或 $p=\infty$ 的情形。
当 $q=\infty$ 时，结论以 $r=p<\infty$ 成立；
当 $p=\infty$ 时，以 $r=q<\infty$ 成立。总有 $r\le\min\{p,q\}$。
\end{remark}

\begin{remark}\label{rem:B-4-8}
在引理~\ref{lem:B-4-6} 的假设下，只要 $r<2$，就有
$q\le2r/(2-r)$。因而有包含 $W^{1,r}\hookrightarrow L^q$。
\end{remark}

\begin{proof}[引理~\ref{lem:B-4-6} 的证明]
设 $\psi:\Omega\to\mathbb R^{2n}$ 是具有紧支集的光滑检验函数，满足
\begin{equation}\label{eq:B-4-4}
  \psi(s,0)\in\mathbb R^n\times\{0\},\qquad
  \partial_t\psi(s,0)\in\{0\}\times\mathbb R^n
\end{equation}
对所有满足 $(s,0)\in\Omega$ 的 $s\in\mathbb R$ 成立。
于是函数 $\phi:=\partial_s\psi-J^T\partial_t\psi$ 属于
$W^{1,p}(\Omega,\mathbb R^{2n})$，满足
$\phi(\Omega\cap\mathbb R)\subset\mathbb R^n\times\{0\}$，并且
\[
  \partial_s\phi+J^T\partial_t\phi
  =\Delta\psi-(\partial_sJ)^T\partial_t\psi
  +(\partial_tJ)^TJ^T\partial_t\psi.
\]
由于在引理的假设下，\eqref{eq:B-4-3} 对紧支于 $\Omega$ 的
$W^{1,p}$ 类检验函数 $\phi$ 仍成立，我们得到
\begin{equation}\label{eq:B-4-5}
  \int_\Omega\langle\Delta\psi,u\rangle
  =-\int_\Omega\langle\partial_s\psi,f\rangle
   -\int_\Omega\langle\partial_t\psi,g\rangle,
\end{equation}
其中
\begin{equation}\label{eq:B-4-6}
  f:=(\partial_tJ)u+\eta,\qquad
  g:=-(\partial_sJ)u-J\eta.
\end{equation}
由假设和 Hölder 不等式，$f,g$ 均属于 $L^r_{\mathrm{loc}}$。
式~\eqref{eq:B-4-5} 表明 $u$ 是
\[
  \Delta u=\partial_sf+\partial_tg
\]
的\BookRecordIndex{T0295}弱解，这里右端应按\BookRecordIndex{T0123}分布意义理解。

先假设 $\Omega\cap\mathbb R=\varnothing$。
由定理~\ref{thm:B-3-2} 得 $u\in W^{1,r}_{\mathrm{loc}}$。
此外，对每个闭包包含在 $\Omega$ 内的有界开集 $\Omega'$，有估计
\begin{align*}
  \|u\|_{W^{1,r}(\Omega')}
  &\le c_1\bigl(\|f\|_{L^r(\Omega)}+\|g\|_{L^r(\Omega)}
       +\|u\|_{L^r(\Omega)}\bigr)\\
  &\le c_1\bigl(\|J\|_{W^{1,p}(\Omega)}\|u\|_{L^q(\Omega)}
       +\|u\|_{L^r(\Omega)}
       +\|J\|_{L^\infty(\Omega)}\|\eta\|_{L^r(\Omega)}
       +\|\eta\|_{L^r(\Omega)}\bigr)\\
  &\le c_2\bigl(1+\|J\|_{W^{1,p}(\Omega)}\bigr)
       \bigl(\|\eta\|_{L^r(\Omega)}+\|u\|_{L^q(\Omega)}\bigr).
\end{align*}
这里 $c_1,c_2$ 依赖于 $p,q,\Omega',\Omega$，而不依赖于 $u,J$。
在这些估计中，我们不失一般性地假设 $\Omega$ 是有界 \BookRecordIndex{T0035}Lipschitz 域，
从而可使用\BookRecordIndex{T0318}索伯列夫嵌入定理。分部积分表明 $u$ 几乎处处满足
\eqref{eq:B-4-2}。这证明了 $\Omega\cap\mathbb R=\varnothing$ 的情形。

为证明一般情形，考虑开集
\[
  \Omega_0:=\{s+it\in\mathbb C\mid s+i|t|\in\Omega\},
\]
并对 $t>0$ 以
\[
  u(s,-t):=Ru(s,t),\qquad f(s,-t):=Rf(s,t),\qquad
  g(s,-t):=-Rg(s,t)
\]
将 $u,f,g$ 延拓至 $\Omega_0$，其中
\[
  R:=\begin{pmatrix}\mathbf 1&0\\0&-\mathbf 1\end{pmatrix}.
\]
这些延拓函数对每个检验函数 $\psi\in C_0^\infty(\Omega_0)$
满足把 $\Omega$ 换成 $\Omega_0$ 后的~\eqref{eq:B-4-5}。
为此，写
\[
  \psi(s,t)=\psi_0(s,t)+\alpha_0(s)\beta_0(t)+\alpha_1(s)\beta_1(t),
\]
其中 $\psi_0$ 满足~\eqref{eq:B-4-4}，且
\[
  \alpha_0(s)\in\{0\}\times\mathbb R^n,\qquad
  \alpha_1(s)\in\mathbb R^n\times\{0\},
\]
\[
  \beta_0(t)=\beta_0(-t),\qquad
  \beta_1(t)=-\beta_1(-t),\qquad
  \beta_0(0)=\partial_t\beta_1(0)=1.
\]
注意 $\alpha_0,\alpha_1$ 由下式确定：
\[
  \psi(s,0)=\psi_0(s,0)+\alpha_0(s),\qquad
  \partial_t\psi(s,0)=\partial_t\psi_0(s,0)+\alpha_1(s).
\]
由 $u,f,g$ 的对称性，对 $i=0,1$ 有
\[
  \int_{\Omega_0}\langle\Delta(\alpha_i\beta_i),u\rangle
  =\int_{\Omega_0}\langle\partial_s(\alpha_i\beta_i),f\rangle
  =\int_{\Omega_0}\langle\partial_t(\alpha_i\beta_i),g\rangle=0.
\]
因此可在~\eqref{eq:B-4-5} 两边把 $\psi$ 换成 $\psi_0$，
而对 $\psi=\psi_0$，该恒等式已经成立。
于是从 $\Omega\cap\mathbb R=\varnothing$ 的情形可知
$u\in W^{1,r}_{\mathrm{loc}}(\Omega_0,\mathbb R^{2n})$，并满足所需估计。

还须证明几乎每个 $s\in\mathbb R\cap\Omega$ 都满足
$u(s,0)\in\mathbb R^n\times\{0\}$。
为此，取一列光滑函数 $v_i:\Omega_0\to\mathbb R^{2n}$，
使其在 $\Omega_0$ 的每个紧子集上以 $W^{1,r}$ 范数收敛到 $u$。
\sourcefootnote{原书将 $v_i$ 的陪域印作 $\mathbb R$，但下式要计算
$Rv_i$ 并与 $u:\Omega_0\to\mathbb R^{2n}$ 比较；据映射类型改为
$\mathbb R^{2n}$。}
于是序列
\[
  u_i(s,t):=\frac12\bigl(v_i(s,t)+Rv_i(s,-t)\bigr)
\]
在 $\Omega_0$ 的每个紧子集上也以 $W^{1,r}$ 范数收敛到 $u$。
此外，对每个 $s\in\mathbb R\cap\Omega$，有
$u_i(s,0)\in\mathbb R^n\times\{0\}$。
由命题~\ref{prop:B-1-21}，$u_i$ 在 $\mathbb R\cap\Omega$
的每个紧区间上以 $L^r$ 范数收敛到 $u$。
由于映射 $s\mapsto u_i(s,0)$ 与投影
$\mathbb R^{2n}\to\{0\}\times\mathbb R^n$ 的复合对每个 $i$ 都为零，
且在 $L^r_{\mathrm{loc}}$ 中收敛到映射
$s\mapsto u(s,0)$ 与同一投影的复合，
可知几乎每个 $s\in\mathbb R\cap\Omega$ 均满足
$u(s,0)\in\mathbb R^n\times\{0\}$。
\sourcefootnote{原书此处极限迹写作 $s\mapsto u(s,t)$；由所用迹命题
及同句中的边界投影，改为 $s\mapsto u(s,0)$。}
\end{proof}

\begin{proposition}\label{prop:B-4-9}
设 $\Omega'\subset\Omega\subset\mathbb H$ 是开集，满足
$\overline{\Omega'}\subset\Omega$\sourcefootnote{原书沿用引理~\ref{lem:B-4-6}
的开集条件。局部\BookRecordIndex{T0408}正则性断言不受影响；但 (ii) 中对整个 $\Omega'$ 的
统一估计，不能仅从未限定有界性的闭包包含关系\BookRecordIndex{T0326}推出，参见前述来源脚注。}，
$\ell$ 为正整数，$p>2$。
则对每个 $c_0>0$，存在 $c>0$，使得如下结论成立。
假设 $J\in W^{\ell,p}(\Omega,\mathbb R^{2n\times2n})$ 满足
$J^2=-\mathbf 1$、\eqref{eq:B-4-1}，且
$\|J\|_{W^{\ell,p}(\Omega)}\le c_0$。
则对每个 $k\in\{0,\ldots,\ell\}$，有以下结论。
\begin{enumerate}[label=\textup{(\roman*)}]
\item 若 $u\in L^p_{\mathrm{loc}}(\Omega,\mathbb R^{2n})$ 和
$\eta\in W^{k,p}_{\mathrm{loc}}(\Omega,\mathbb R^{2n})$
使~\eqref{eq:B-4-3} 成立，则
$u\in W^{k+1,p}_{\mathrm{loc}}(\Omega,\mathbb R^{2n})$，
且 $u$ 几乎处处满足~\eqref{eq:B-4-2}。
\item 若 $u\in W^{k+1,p}_{\mathrm{loc}}(\Omega)$ 满足
$u(\Omega\cap\mathbb R)\subset\mathbb R^n\times\{0\}$，则
\[
  \|u\|_{W^{k+1,p}(\Omega')}
  \le c\bigl(\|\partial_su+J\partial_tu\|_{W^{k,p}(\Omega)}
          +\|u\|_{W^{k,p}(\Omega)}\bigr).
\]
\end{enumerate}
\end{proposition}

\begin{proof}
证明分三步。

\emph{步骤 1：$k=0,\ell=1$ 时结论成立。}
由引理~\ref{lem:B-4-6} 和\BookRecordIndex{T0318}索伯列夫嵌入定理，
\[
  u\in L^q_{\mathrm{loc}},\quad
  \frac p{p-1}<q<\frac{2p}{p-2}
  \quad\Longrightarrow\quad
  u\in W^{1,r}_{\mathrm{loc}}\subset L^{q'}_{\mathrm{loc}},
\]
其中
\[
  r:=\frac{pq}{p+q}<2,\qquad q':=\frac{2r}{2-r}>q.
\]
现按引理~\ref{lem:B-4-4} 取 $q_j,r_j$，$j=0,\ldots,m$。
由于 $q_0\le p$，有 $u\in L^{q_0}_{\mathrm{loc}}$；
归纳得到对 $j=0,\ldots,m$ 均有 $u\in W^{1,r_j}_{\mathrm{loc}}$。
取 $j=m$，则 $q_m>2p/(p-2)$，因而
\[
  r_m=\frac{pq_m}{p+q_m}>2,
\]
所以 $u$ 连续。再取 $q=\infty,r=p$ 应用引理~\ref{lem:B-4-6}
（见注~\ref{rem:B-4-7}），得到 $u\in W^{1,p}_{\mathrm{loc}}$。
这证明了 $k=0,\ell=1$ 时的 (i)。

为证明 (ii)，取 $\Omega$ 内有界的开 \BookRecordIndex{T0035}Lipschitz 域
$\Omega_0,\ldots,\Omega_m$，满足
\[
  \overline{\Omega'}\subset\operatorname{int}_{\mathbb H}(\Omega_m),
  \qquad \overline{\Omega_j}\subset
  \operatorname{int}_{\mathbb H}(\Omega_{j-1}),
  \qquad \overline{\Omega_0}\subset\Omega.
\]
则对 $j=1,\ldots,m$，由引理~\ref{lem:B-4-6} 得
\begin{align*}
  \|u\|_{W^{1,r_j}(\Omega_j)}
  &\le c_1\bigl(\|\partial_su+J\partial_tu\|_{L^{r_j}(\Omega_{j-1})}
       +\|u\|_{L^{q_j}(\Omega_{j-1})}\bigr)\\
  &\le c_2\bigl(\|\partial_su+J\partial_tu\|_{L^p(\Omega)}
       +\|u\|_{W^{1,r_{j-1}}(\Omega_{j-1})}\bigr)\\
  &\le c_3\bigl(\|\partial_su+J\partial_tu\|_{L^p(\Omega)}
       +\|u\|_{W^{1,r_0}(\Omega_0)}\bigr)\\
  &\le c_4\bigl(\|\partial_su+J\partial_tu\|_{L^p(\Omega)}
       +\|u\|_{L^p(\Omega)}\bigr).
\end{align*}
第二个不等式由\BookRecordIndex{T0318}索伯列夫嵌入定理\BookRecordIndex{T0326}推出，第三个由归纳\BookRecordIndex{T0326}推出；
最后一个不等式用到了 $q_0\le p$。
取 $j=m$，有 $r_m>2$，故由定理~\ref{thm:B-1-11}
得到 $u$ 的 $L^\infty$ 范数估计。然后再次以 $q=\infty,r=p$
使用引理~\ref{lem:B-4-6} 的估计。

\emph{步骤 2：$k=\ell=1$ 时结论成立。}
取 $q,r>1$，满足 $1/p+1/q=1/r$。则
\begin{equation}\label{eq:B-4-7}
  u\in W^{1,q}_{\mathrm{loc}}
  \quad\Longrightarrow\quad u\in W^{2,r}_{\mathrm{loc}},
\end{equation}
且 $u$ 满足相应估计。为证明此点，注意函数
\[
  u':=\partial_su\in L^q_{\mathrm{loc}},\qquad
  \eta':=\partial_s\eta-(\partial_sJ)\partial_tu\in L^r_{\mathrm{loc}}
\]
对每个 $\phi\in C_0^\infty(\Omega,\mathbb R^{2n})$
满足~\eqref{eq:B-4-3}。故由引理~\ref{lem:B-4-6}，
$u'\in W^{1,r}_{\mathrm{loc}}$。
又由引理~\ref{lem:B-4-5}，
$\partial_tu=J(\partial_su-\eta)\in W^{1,r}_{\mathrm{loc}}$，
从而 $u\in W^{2,r}_{\mathrm{loc}}$。

为证明 $k=1$ 时的断言，按引理~\ref{lem:B-4-4} 取序列
$q_j,r_j$，并如步骤 1 那样归纳使用~\eqref{eq:B-4-7}，
得到对每个 $j$ 均有 $u\in W^{2,r_j}_{\mathrm{loc}}$。
取 $j=m$，$u$ 连续可微；再取 $q=\infty$ 使用~\eqref{eq:B-4-7}，
便有 $u\in W^{2,p}_{\mathrm{loc}}$。
这证明了 $k=\ell=1$ 时的 (i)。
为证明 (ii)，完全按步骤 1 论证；此处细节留给读者。

\emph{步骤 3：一般情形。}
归纳假设命题对某个 $k\ge1$ 成立，且 $J\in W^{k+1,p}$。
设 $\eta\in W^{k+1,p}_{\mathrm{loc}}$。
将归纳假设应用于步骤 2 中的 $u',\eta'$，得到
$\partial_su=u'$ 和 $\partial_tu=J(\partial_su-\eta)$
均为 $W^{k+1,p}_{\mathrm{loc}}$ 类，并满足形如
\begin{align*}
  \|u\|_{W^{k+2,p}(\Omega')}
  &\le c_1\bigl(\|u\|_{W^{k+1,p}(\Omega')}
       +\|u'\|_{W^{k+1,p}(\Omega')}
       +\|\eta\|_{W^{k+1,p}(\Omega')}\bigr)\\
  &\le c_2\bigl(\|\eta\|_{W^{k+1,p}(\Omega)}
       +\|u\|_{W^{k,p}(\Omega)}
       +\|\eta'\|_{W^{k,p}(\Omega)}
       +\|u'\|_{W^{k,p}(\Omega)}\bigr)\\
  &\le c_3\bigl(\|\eta\|_{W^{k+1,p}(\Omega)}
       +\|u\|_{W^{k+1,p}(\Omega)}\bigr)
\end{align*}
的估计。第一个不等式由命题~\ref{prop:B-1-19} \BookRecordIndex{T0326}推出
（因为 $\partial_tu=J(u'-\eta)$ 且 $(k+1)p>2$）；
第二个不等式由分别应用于 $u,u'$ 的归纳假设\BookRecordIndex{T0326}推出；
最后一个不等式由命题~\ref{prop:B-1-19} \BookRecordIndex{T0326}推出
（因为 $\eta'=\partial_s\eta-(\partial_sJ)\partial_tu$ 且 $kp>2$）。
命题~\ref{prop:B-4-9} 得证。
\end{proof}

\begin{proof}[定理~\ref{thm:B-4-1} 的证明]
只需在 $\Sigma$ 的局部全纯坐标和 $M$ 的局部坐标中证明结论。
若 $J$ 是 $C^\ell$ 类\BookRecordIndex{T0088}殆复结构，则 $M$ 上存在 $C^\ell$ 局部坐标卡，
它把 $L$ 与 $\mathbb R^n\times\{0\}$ 识别，且沿 $L$ 把 $J$ 与 $J_0$
识别（见下面的习题~\ref{ex:B-4-10}）。
用 $C^\ell$ 坐标卡推前 $J$，得到 $\mathbb R^{2n}$ 上的
$C^{\ell-1}$ 类\BookRecordIndex{T0088}殆复结构。因此可假设 $\Omega\subset\mathbb H$ 是开集，
$u:\Omega\to\mathbb R^{2n}$ 是边值问题
\[
  \partial_su+J(u)\partial_tu=0,\qquad
  u(\Omega\cap\mathbb R)\subset\mathbb R^n\times\{0\}
\]
的解，其中 $J:\mathbb R^{2n}\to\mathbb R^{2n\times2n}$
是 $C^{\ell-1}$ 类\BookRecordIndex{T0088}殆复结构，并满足
\begin{equation}\label{eq:B-4-8}
  x\in\mathbb R^n\times\{0\}\quad\Longrightarrow\quad J(x)=J_0.
\end{equation}
先取 $k=1$。则 $u\in W^{1,p}_{\mathrm{loc}}(\Omega,\mathbb R^{2n})$
满足命题~\ref{prop:B-4-9} 的条件，其中 $k=1$、$\eta=0$，
并用 $J\circ u\in W^{1,p}_{\mathrm{loc}}(\Omega,\mathbb R^{2n\times2n})$
代替 $J$。故 $u\in W^{2,p}_{\mathrm{loc}}$。
继续归纳，在有限 $\ell$ 的情形依次得到
$u\in W^{k+1,p}_{\mathrm{loc}}$，$k=1,\ldots,\ell-1$。
因而 $u\in W^{\ell,p}_{\mathrm{loc}}$；若 $\ell=\infty$，
则对每个有限阶数重复这一论证，得到 $u$ 光滑。
\sourcefootnote{原书在第一次提升后立即写 $J\circ u\in W^{2,p}_{\mathrm{loc}}$；
若仅有 $\ell=2$，推前的 $J$ 只有 $C^1$，该二阶复合\BookRecordIndex{T0408}正则性并不由此\BookRecordIndex{T0326}推出，
而且此时也无需再作下一步提升。这里按实际所需的归纳阶数表述。}
定理~\ref{thm:B-4-1} 得证。
\end{proof}

\begin{proof}[定理~\ref{thm:B-4-2} 的证明]
由于 $p>2$ 时包含 $W^{1,p}\hookrightarrow C$ 是紧的，
可假设 $u_\nu$ 在 $\Sigma$ 的紧子集上一致收敛到连续映射
$u:\Sigma\to M$。\sourcefootnote{原书仅说“一致收敛”，
但假设只在逐个紧集上给出 $W^{1,p}$ 一致界；当 $\Sigma$ 非紧时，
此处由局部紧嵌入与对角子序列得到的是\BookRecordIndex{T0416}逐紧一致收敛。}
故由下面的习题~\ref{ex:B-4-10}，可假设
$u_\nu:\Omega\to\mathbb R^{2n}$ 是一列 $W^{1,p}_{\mathrm{loc}}$ 类解，
满足一列边值问题；此处将原书的边界子空间补为
$\mathbb R^n\times\{0\}$。\sourcefootnote{原书将它印作
$\mathbb R^{2n}\times\{0\}$，与 $u_\nu$ 的靶 $\mathbb R^{2n}$
类型不合，且与紧邻的边界条件~\eqref{eq:B-4-8} 矛盾。}
\[
  \partial_su_\nu+J_\nu(u_\nu)\partial_tu_\nu=0,\qquad
  u_\nu(\Omega\cap\mathbb R)\subset\mathbb R^n\times\{0\},
\]
其中 $J_\nu\in C^{\ell-1}(\mathbb R^{2n},\mathbb R^{2n\times2n})$
在 $C^{\ell-1}$ 拓扑下收敛到 $J$，并满足~\eqref{eq:B-4-8}。
由命题~\ref{prop:B-4-9} 归纳可知，$u_\nu\in W^{\ell,p}_{\mathrm{loc}}$，
且对每个具有光滑边界的紧子集 $Q\subset\Omega$，有
\[
  \sup_\nu\|u_\nu\|_{W^{\ell,p}(Q)}<\infty.
\]
由定理~\ref{thm:B-1-11}，包含
$W^{\ell,p}(Q)\hookrightarrow C^{\ell-1}(Q)$ 是紧的。
因此在 $\Sigma$ 的任意给定紧子集上，存在一个在 $C^{\ell-1}$ 拓扑下
收敛的子序列。取 $\Sigma$ 的一列穷竭紧子集，再取对角子序列，
便得定理~\ref{thm:B-4-2}。
\end{proof}

\begin{exercise}\label{ex:B-4-10}
设 $J\in\mathcal J^\ell(M,L)$。证明：对 $L$ 中每个点，
存在邻域 $U$ 和一个 $C^\ell$ 坐标卡 $\psi:U\to\mathbb R^{2n}$，使得
\[
  \psi(U\cap L)=\psi(U)\cap(\mathbb R^n\times\{0\}),\qquad
  J_0d\psi(x)=d\psi(x)J(x)
\]
对 $x\in U\cap L$ 成立。若
$J_\nu\in\mathcal J^\ell(M,L)$ 在 $C^\ell$ 拓扑下收敛到 $J$，
证明可选取 $\psi_\nu$，使其在 $C^\ell$ 拓扑下收敛到 $\psi$。
提示：设 $\phi:V\to\mathbb R^n$ 是 $L$ 的开集 $V$ 上的坐标卡
（它可延拓到 $V$ 的闭包），令
\[
  U_J:=\{\exp_x(J(x)\xi)\mid x\in V,\ \xi\in T_xL,\ |\xi|<\delta\},
\]
并定义 $\psi_J:U_J\to\mathbb R^{2n}$ 为
\[
  \psi_J(\exp_x(J(x)\xi)) := (\phi(x),d\phi(x)\xi).
\]
\end{exercise}

\begin{exercise}\label{ex:B-4-11}
设 $\Omega\subset\mathbb H$ 是上半平面的开子集。
假设 $J\in C^\infty_{\mathrm{loc}}(\Omega,\mathbb R^{2n\times2n})$
满足 $J^2=-\mathbf 1$ 和~\eqref{eq:B-4-1}。
证明：命题~\ref{prop:B-4-9} 中的假设 $p>2$ 可以去掉。
\end{exercise}

\begin{exercise}\label{ex:B-4-12}
当 $J$ 可积且 $\Sigma$ 为闭曲面时，找出你所能给出的
定理~\ref{thm:B-4-1} 和定理~\ref{thm:B-4-2} 的最简证明。
\end{exercise}
